\begin{table}[p]
\centering
\caption{Existing-artifact 2 by 2 candidate-bank and locking-rule comparison for matched AMP-AD development conditions.}
\label{tab:supp19}
\scriptsize
\setlength{\tabcolsep}{2.6pt}
\renewcommand{\arraystretch}{1.08}
\resizebox{\linewidth}{!}{%
\begin{tabular}{llrrrrrrr}
\toprule
Candidate bank & Locking rule & $n$ & Total target deviation & Mean selected $n$ & Mean score & Mean regret & Maximum regret & Mask changes \\
\midrule
Original-objective bank & Original top-three & 24 & 216 & 24.12 & 0.86475 & 0.00261 & 0.02111 & 10 \\
Original-objective bank & Regret-constrained & 24 & 227 & 24.58 & 0.86676 & 0.00060 & 0.00729 & 10 \\
Updated-objective bank & Original top-three & 24 & 135 & 20.75 & 0.86329 & 0.00351 & 0.01525 & 8 \\
Updated-objective bank & Regret-constrained & 24 & 137 & 20.83 & 0.86572 & 0.00108 & 0.00835 & 8 \\
\bottomrule
\end{tabular}}
\par\smallskip
\parbox{0.98\linewidth}{\scriptsize\textit{Note.} No candidate generation or held-out evaluation was rerun. All 24 original-objective banks and all 24 updated-objective banks were complete after recovery of the 30 Rush run-best feature lists from the authoritative Drive source folder. ``Mask changes'' is the within-bank number of conditions for which the two locking rules selected different masks. The regret-constrained rule improved the configured development-CV score-gap summaries in both bank sets but did not universally improve target fidelity.}
\end{table}
